\chapter{正規表達式}
\label{ch:regex}

到目前為止，你已經用 \texttt{in}、\texttt{.replace()}、\texttt{.split()} 這些字串方法處理過
不少「文字裡有沒有出現某個東西、要不要換掉某段文字」的問題。這些方法有一個共同的前提：你得
\textbf{確切知道}要找的字串長什麼樣子。

但現實裡很多時候你要找的其實是一種\textbf{格式}：
你可能不知道一支電話號碼的確切數字是什麼，但看到「三個數字、一個連字號、三個數字、一個連字
號、四個數字」這種形狀，你一眼就認得出那是電話號碼；一個 email 地址也是一樣，不管帳號跟網域
是什麼，「一段文字、一個 \texttt{@}、再一段文字、一個點」這個形狀是固定的。正規表達式
（regular expression，常簡稱 regex）就是讓你把這種形狀寫給 Python 看的工具。
\index{正規表達式}

這一章分幾段推進：

\begin{itemize}
\item 先從你已經會的字串方法出發，看看「找格式」這件事跟「找固定字串」有什麼不
一樣，順便學會 \texttt{re} 模組最基本的用法。
\item 接著把「形狀」拆解成幾個最常用的積木（數字、
文字、空白、重複次數），組出電話號碼、email 這類實際會用到的 pattern。
\item 再來學怎麼一次把一段
文字裡所有符合的地方都找出來、怎麼確認一整段文字完全符合某種格式（也就是「驗證」）、怎麼用
regex 做搜尋取代跟切字串。
\end{itemize}

要先說在前面：正規表達式是一個很深的主題，光是語法規則就能寫成一
整本書，但這一章的目標不是讓你精通它，只要看到常見的 pattern 看得懂在比對什麼、需要的時候
自己也能寫出簡單的 pattern（數字、email 這種格式），這一章的任務就完成了，更進階的寫法留給
你以後真的用到時再自己查文件。

\section{為什麼需要正規表達式、\texttt{re} 模組怎麼用}
\label{sec:regex-basics}

假設你在整理一批客服留言，想知道哪些留言裡有留電話號碼、哪些沒有。用你已經會的工具，你會
怎麼寫？如果你知道確切的電話號碼是什麼，\texttt{in} 可以幫你判斷這段文字有沒有出現；但你
不可能事先知道每個客戶會留下哪個號碼，你真正想問的其實是「這裡有沒有出現\textbf{一串數字}」，
而不是「這裡有沒有出現某個特定的字串」。

這正是字串方法做不到、但正規表達式擅長的事：它讓你
可以描述「找長這種樣子的東西」，不只是「找這個字」。

要在 Python 裡用正規表達式，第一步是 \texttt{import re}。\texttt{re} 是 Python 內建的正規
表達式模組，不用另外安裝。最常用的入口是 \texttt{re.search(pattern, text)}：如果 \texttt{text}
裡有一段符合 \texttt{pattern} 這個「形狀」的文字，它會回傳一個 Match 物件；找不到的話，回傳
\texttt{None}。拿到 Match 物件之後，呼叫它的 \texttt{.group()} 方法，就能取出實際比對到的那段
文字。
\index{re 模組}
\index{re.search()}
\index{Match 物件}

\begin{code}{python}{Matching a pattern and extracting the matched text}
import re

match = re.search(pattern, text)
if match:
    print(match.group())
else:
    print("Not found")
\end{code}

這裡 \texttt{pattern} 通常會寫成 \texttt{r'...'} 這種前面多一個 \texttt{r} 的字串，叫 raw
string。正規表達式常常要用到反斜線（像等一下會教的 \texttt{\textbackslash d}），如果不加
這個 \texttt{r}，反斜線可能被 Python 字串本身的跳脫規則誤解成別的意思。之後每一段 pattern，
這本書都會統一用 \texttt{r'...'} 這種寫法，你也照著做就好，不用另外深究背後的規則。
\index{raw string}

比對「一個數字」的 pattern 寫作 \texttt{\textbackslash d}，先只用這一個最基本的元件，其他
常用的元件留到下一節再教。

\begin{example}{python}{(check-digit.py) Checking whether a message contains a digit}
import re

message = "Hello, my phone number is 0912345678, please call back soon."

match = re.search(r'\d', message)
if match:
    print("This message contains a digit, the first one is:", match.group())
else:
    print("This message has no digits")
\end{example}

\iosample{}{This message contains a digit, the first one is: 0 \\}

如果同一個 pattern 之後還要重複用好幾次，你也可以先用 \texttt{re.compile()} 把它變成一個
regex 物件、貼給一個名字，例如 \texttt{pattern = re.compile(r'\textbackslash d')}，之後
每次呼叫 \texttt{pattern.search(text)} 就好。這裡的 \texttt{pattern} 這個名字牌，跟你前面
幾章看過的物件用法完全一樣，只是這次貼的對象是一個 regex 物件而已。

\index{re.compile()}
還記得第 \ref{ch:files-exceptions} 章教的「先做、錯了再處理」（EAFP）嗎？這裡看起來不太一樣：\texttt{search()} 比對
不到的時候，是老實回傳 \texttt{None}，不是丟例外，所以慣用寫法是先用 \texttt{if match:} 檢查
再用，乍看之下比較像「先檢查」（LBYL）。

這其實仍然算是 EAFP 的精神，只是同一個原則的另一面：
EAFP 真正在乎的是操作本身可能失敗時，要用能反映這件事的寫法。\texttt{search()} 的設計本來
就是老實回答「有」或「沒有」，不是「可能出錯會丟例外的操作」（不像開檔案、字典查詢那樣），所以
檢查 \texttt{if match:} 才是正確的寫法，不需要每個地方都硬套 \texttt{try/except}。

如果忘記檢查、直接對比對結果呼叫 \texttt{.group()}，會發生什麼事？

\begin{example}{python}{(check-digit-broken.py) Calling .group() without checking first}
import re

message = "Thank you for your message, a representative will get back to you soon."
match = re.search(r'\d', message)
print(match.group())
\end{example}

這則留言裡沒有任何數字，\texttt{match} 會是 \texttt{None}，執行到最後一行時會看到：

\begin{verbatim}
Traceback (most recent call last):
  File "check-digit-broken.py", line 4, in <module>
    print(match.group())
          ^^^^^^^^^^^
AttributeError: 'NoneType' object has no attribute 'group'
\end{verbatim}

由下往上讀：最後一行 \texttt{AttributeError: 'NoneType' object has no attribute 'group'}
是說「\texttt{None} 這種東西身上沒有 \texttt{group} 這個方法可以呼叫」，因為 \texttt{match}
根本不是 Match 物件，是 \texttt{search()} 老實告訴你「沒找到」的 \texttt{None}。修法很直接：
呼叫 \texttt{.group()} 之前，一定要先用 \texttt{if match:} 確認真的有比對到，這也是為什麼
本節一開始的片語會把 \texttt{if match:} 當成固定樣板的一部分，不是可有可無的裝飾。

\begin{drills}[label=drills:regex-basics]
\item 用 \texttt{re.search(r'\textbackslash d', "Customer ID A007")} 這行程式，你覺得會不會
比對到？如果會，比對到的是哪個字元？
\item 用 \texttt{re.search(r'\textbackslash d', "A sentence without any digits")}，你覺得 \texttt{match}
最後會是 \texttt{None} 還是一個 Match 物件？如果接著不檢查、直接寫
\texttt{match.group()}，會發生什麼事？
\end{drills}

\section{常用 pattern 積木：\texttt{\textbackslash d}、\texttt{\textbackslash w}、
\texttt{\textbackslash s} 與量詞}
\label{sec:regex-patterns}

上一節只用了「比對一個數字」這個最小的積木，但實際情境裡，你通常要描述的是一整段有規律的
格式，例如「三個數字、一個連字號、三個數字、一個連字號、四個數字」。要做到這件事，你需要
兩種積木：一種是「這個位置該是什麼樣的字元」，另一種是「這個位置要重複幾次」。

最常用的三個字元類別：

\begin{itemize}
\item \texttt{\textbackslash d}：一個數字（0 到 9 之間任何一個）。
\item \texttt{\textbackslash w}：一個「文字字元」，包含英文字母、數字、底線。
\item \texttt{\textbackslash s}：一個空白字元，包含空格、換行、tab。
\end{itemize}
\index{字元類別 d, w, s}

拿標點符號試試就能感受到 \texttt{\textbackslash w} 跟你想像的「文字」不完全一樣：
\texttt{re.search(r'\textbackslash w', "！")} 會是 \texttt{None}，因為驚嘆號不算英文
字母、數字或底線；而 \texttt{re.search(r'\textbackslash s', "a b").group()} 會找到
\texttt{"a"} 跟 \texttt{"b"} 中間那個空格。

光有字元類別還不夠，你還需要說「要重複幾次」，這就是量詞的工作：

\begin{itemize}
\item \texttt{*}：前面那個元件出現零次以上（可以完全不出現）。
\item \texttt{+}：前面那個元件出現一次以上（至少要有一個）。
\item \texttt{?}：前面那個元件可有可無（出現零次或一次）。
\item \texttt{\{n\}}：前面那個元件精確出現 \texttt{n} 次。
\item \texttt{\{n,m\}}：前面那個元件出現 \texttt{n} 到 \texttt{m} 次（含頭尾兩端）。
\end{itemize}
\index{量詞}

把字元類別跟量詞組合起來，就能寫出「三個數字、一個連字號、三個數字、一個連字號、四個數字」
這種格式：\texttt{\textbackslash d\{3\}-\textbackslash d\{3\}-\textbackslash d\{4\}}。

\begin{code}{python}{Combining character classes with quantifiers}
import re

pattern = r'\d{3}-\d{3}-\d{4}'
match = re.search(pattern, text)
if match:
    print(match.group())
\end{code}

\begin{example}{python}{(check-phone-format.py) Checking whether a message has a phone number in the required format}
import re

message = "Hello, my phone number is 091-234-5678, please call back soon."

match = re.search(r'\d{3}-\d{3}-\d{4}', message)
if match:
    print("Found a phone number matching the format:", match.group())
else:
    print("No phone number matching the format was found")
\end{example}

\iosample{}{Found a phone number matching the format: 091-234-5678 \\}

如果客戶留的號碼格式不太一樣，例如少了一個連字號、或分段方式不同，就比對不到。把
\texttt{message} 換成 \texttt{"Hello, my phone number is 0912-345-678, please call back soon."} 再跑一次，
會印出「No phone number matching the format was found」，因為 \texttt{4-3-3} 這種分段方式跟
\texttt{\textbackslash d\{3\}-\textbackslash d\{3\}-\textbackslash d\{4\}} 描述的
\texttt{3-3-4} 形狀對不上。

\keypoint{字元類別跟量詞要記住的重點是「規律的東西才需要 regex」，不是要你把每個符號背下
來。如果你要找的字串本身就是固定的（例如一定是 \texttt{"客服專線"} 這四個字），直接用
\texttt{in} 或 \texttt{.replace()} 更簡單、也更好懂，不用每個地方都想著要動用正規表達式。}

\begin{drills}[label=drills:regex-patterns]
\item 下面三個字串，哪些符合 \texttt{\textbackslash d\{3\}-\textbackslash d\{3\}-%
\textbackslash d\{4\}} 這個 pattern？為什麼：
(a) \texttt{"091-234-5678"}
(b) \texttt{"0912-345-678"}
(c) \texttt{"12-345-6789"}
\item 如果要用同樣的精神，寫一個「兩碼-兩碼-兩碼-兩碼」格式的 pattern（例如
\texttt{"12-34-56-78"}），你會怎麼寫？
\item 句子 \texttt{"Customer service line: 091-234-5678, service hours 9:00-18:00"} 用
\texttt{re.search(r'\textbackslash d\{3\}-\textbackslash d\{3\}-\textbackslash d\{4\}',
sentence)} 去找，比對到的會是哪一段文字？
\item 用 \texttt{\{n,m\}} 這個新學到的量詞，寫一個「密碼長度要在 6 到 12 個文字字元之間」
的 pattern（提示：文字字元就是上面教的 \texttt{\textbackslash w}）。
\end{drills}

\section{\texttt{findall()}：一次找出所有符合的地方}
\label{sec:regex-findall}

前面兩節的 \texttt{re.search()} 只會告訴你「有沒有找到、第一個找到的是哪段」，但很多實際
任務要的是這段文字裡\textbf{所有}符合的地方，例如從一份聯絡資訊清單裡，一次把所有 email
地址都抓出來。這種時候要用 \texttt{re.findall(pattern, text)}：它會把 \texttt{text} 裡所有
符合 \texttt{pattern} 的地方，通通收集成一個 list 回傳給你。
\index{re.findall()}

在動手之前，先組出「email 地址」這個 pattern。一個 email 大致長這樣：帳號部分（英文字母、
數字，可能夾雜幾個符號）、一個 \texttt{@}、網域部分（英文字母、數字、點）、最後接一段純
英文字母。把這幾段用上一節的積木拼起來，會得到
\texttt{[a-zA-Z0-9.\_\%+-]+@[a-zA-Z0-9.-]+\textbackslash.[a-zA-Z]+}。

中括號
\index{字元集合}
\texttt{[...]} 是「這個位置可以是括號裡列出的任何一個字元」，\texttt{+} 表示「這樣的字元
至少出現一次」，\texttt{\textbackslash.} 則是用反斜線跳脫，表示「真正的一個點，不是別的
意思」。這個 pattern 已經比前面幾節複雜不少，你不需要自己從零想出這種寫法，記得它長什麼
樣子、看得懂在比對什麼，需要的時候查文件或直接沿用，就已經達到這一章的目標。

「點要記得跳脫」是初學正規表達式很容易忽略的一個細節，因為 \texttt{.} 這個符號在正規表達式
裡有自己的意思：它代表「任何一個字元都可以」，不是專指「一個真正的點」。如果忘記在點前面加
反斜線，pattern 會比你以為的鬆很多：

\begin{example}{python}{(dot-escape.py) Forgetting to escape the dot makes a pattern too loose}
import re

samples = ["3.14", "3x14", "3-14"]

pattern_loose = r'\d.\d'    # 這裡的 . 沒有跳脫，所以可以比對任何字元
pattern_strict = r'\d\.\d'  # 這裡的 \. 才會真正比對到一個點

for s in samples:
    loose = re.search(pattern_loose, s)
    strict = re.search(pattern_strict, s)
    print(s, "loose:", loose.group() if loose else None,
          "strict:", strict.group() if strict else None)
\end{example}

\iosample{}{3.14 loose: 3.1 strict: 3.1 \\
3x14 loose: 3x1 strict: None \\
3-14 loose: 3-1 strict: None \\}

沒有跳脫的 \texttt{pattern\_loose} 對 \texttt{"3x14"}、\texttt{"3-14"} 這兩個字串也照樣
比對成功，因為 \texttt{.} 把中間那個 \texttt{x} 或 \texttt{-} 都當成「隨便一個字元」接受
了，並非只認得真正的點。加上反斜線的 \texttt{pattern\_strict} 才正確地只認得真正的一個點。
本節開頭 email pattern 裡的 \texttt{\textbackslash.}，用的就是這個跳脫寫法。

\begin{code}{python}{Finding all text that matches a pattern at once}
import re

results = re.findall(pattern, text)
print(results)
\end{code}

\begin{example}{python}{(find-emails.py) Finding all email addresses in a block of contact info}
import re

contacts = """Alice - alice.wang@example.com
Bob - bob99@mail-service.net
Cindy hasn't left an email yet, just wrote a note"""

pattern = r'[a-zA-Z0-9._%+-]+@[a-zA-Z0-9.-]+\.[a-zA-Z]+'
emails = re.findall(pattern, contacts)
print(emails)

# 把 pattern 用括號包起來，findall() 就只會回傳括號裡比對到的那一段
username_pattern = r'([a-zA-Z0-9._%+-]+)@[a-zA-Z0-9.-]+\.[a-zA-Z]+'
usernames = re.findall(username_pattern, contacts)
print(usernames)
\end{example}

\iosample{}{['alice.wang@example.com', 'bob99@mail-service.net'] \\{}
['alice.wang', 'bob99'] \\}

留意最後兩行的差別：第一次的 pattern 沒有括號，\texttt{findall()} 回傳的是完整比對到的
email；第二次把帳號的部分用括號 \texttt{()} 包起來，\texttt{findall()} 就只回傳括號裡比對
到的那一小段。這其實跟第一節 \texttt{match.group()} 是同一件事的延伸：括號告訴 Python
「這一段我特別在意，之後只要把這段還我就好」。這個用法叫「分組」，正式的規則（多個括號怎麼
編號、怎麼分別取出）還有更多細節，這裡先知道加了括號、拿到的東西會變就夠用。
\index{分組}

如果一個 pattern 裡有好幾組括號，光靠「第一個括號、第二個括號」這種順序去記，讀起來會很
吃力。這種時候可以把括號寫成\textbf{命名分組}，語法是 \texttt{(?P<名字>...)}，之後就能用
這個名字取出對應的那一段，而不用去數第幾個括號：
\index{命名分組}

\begin{example}{python}{(named-group-phone.py) Using named groups to pull apart a phone number}
import re

text = "Taylor's phone number is 091-234-5678"
pattern = r'(?P<area>\d{3})-(?P<mid>\d{3})-(?P<last>\d{4})'

match = re.search(pattern, text)
if match:
    print(match.group("area"))
    print(match.group("mid"))
    print(match.group("last"))
\end{example}

\iosample{}{091 \\
234 \\
5678 \\}

命名分組的好處是可讀性：一段複雜的 pattern 裡如果有好幾個括號，用名字取比用數字數第幾個
括號清楚很多。這裡先知道有這個寫法即可，之後真的碰到分組較多、需要清楚標示每一段代表什麼的
情境，再回頭查文件細節。

\section{驗證格式：用 \texttt{fullmatch()} 檢查整段字串}
\label{sec:regex-fullmatch}

前面幾節的 \texttt{search()}、\texttt{findall()} 有一個共同的特性：它們只在乎文字裡「有
沒有一段」符合 pattern，不在乎符合的那一段前後還有沒有其他文字。這對「找出留言裡的電話
號碼」這種任務很合適，但如果你的任務換成「驗證使用者填的這整個欄位是不是一支合法的電話
號碼」，\texttt{search()} 反而會讓你上當：

\begin{example}{python}{(search-trap.py) search() is not suitable for validation}
import re

pattern = r'\d{3}-\d{3}-\d{4}'
text = "091-234-5678abc"

match = re.search(pattern, text)
print("search() found:", match.group())
\end{example}

\iosample{}{search() found: 091-234-5678 \\}

這個輸入其實不是一支合法的電話號碼（後面多了 \texttt{"abc"}），但 \texttt{search()} 只顧
找出「有沒有一段符合」，找到 \texttt{"091-234-5678"} 這一段就滿意收工，完全不管後面還跟著
什麼。要正確做「驗證」，你需要的是確認\textbf{整段字串從頭到尾}都符合 pattern，不能有多餘
的部分，也不能少。這正是 \texttt{re.fullmatch(pattern, text)} 的工作：只有整個 \texttt{text}
完全符合 \texttt{pattern}，才會回傳 Match 物件，只要頭尾任何一邊多了或少了字元，都回傳
\texttt{None}。
\index{re.fullmatch()}

\begin{code}{python}{Confirming a whole string fully matches a pattern}
import re

result = re.fullmatch(pattern, text)
if result:
    print("Format is valid")
else:
    print("Format is invalid")
\end{code}

\begin{example}{python}{(validate-phone.py) Validating whether a phone field is in a legal format}
import re

def is_valid_phone(text):
    pattern = r'\d{3}-\d{3}-\d{4}'
    return re.fullmatch(pattern, text) is not None

print(is_valid_phone("091-234-5678"))
print(is_valid_phone("091-234-5678 "))
print(is_valid_phone("my phone number is 091-234-5678"))
\end{example}

\iosample{}{True \\
False \\
False \\}

三個輸入裡，只有第一個從頭到尾完全符合「三碼-三碼-四碼」的格式。第二個結尾多了一個空格、
第三個開頭多了一段文字，用 \texttt{search()} 都還是找得到，但用 \texttt{fullmatch()} 都會
正確判定「不合法」，因為它們符合的都只是字串中間夾著的一段，整段字串本身並不完全符合格式。

\keypoint{\texttt{search()} 跟 \texttt{fullmatch()} 很容易搞混，因為兩者的呼叫方式幾乎
一樣。判斷用哪一個的標準很簡單：如果任務是「這段文字裡有沒有出現某種格式」，用
\texttt{search()}；如果任務是「使用者填的這整個欄位，是不是一個合法的格式」，用
\texttt{fullmatch()}。用錯的後果是實際的資安風險，例如驗證信箱格式時如果誤用
\texttt{search()}，使用者填 \texttt{"attacker@evil.com; DROP TABLE users"} 這種夾帶額外
內容的字串也會被判定成合法，因為裡面確實藏著一段符合 email pattern 的文字。}

同樣的做法拿來驗證 email 也一樣：

\begin{example}{python}{(validate-email.py) Validating whether an email field is in a legal format}
import re

email_pattern = r'[a-zA-Z0-9._%+-]+@[a-zA-Z0-9.-]+\.[a-zA-Z]+'

def is_valid_email(text):
    return re.fullmatch(email_pattern, text) is not None

for candidate in ["alice.wang@example.com", "alice.wang@example.com please inquire",
                  "alice@", "a@b.tw"]:
    print(candidate, "->", is_valid_email(candidate))
\end{example}

\iosample{}{alice.wang@example.com -> True \\
alice.wang@example.com please inquire -> False \\
alice@ -> False \\
a@b.tw -> True \\}

留意 \texttt{"alice@"} 這個輸入：帳號跟 \texttt{@} 都在，但少了網域部分，不符合 pattern
要求的「\texttt{@} 後面要有網域、一個點、再一段英文字母」，所以正確判定成不合法。這正是
regex 驗證的價值：不用一項一項手動檢查「有沒有 \texttt{@}」「有沒有點」，一個 pattern 就
把整套格式規則講清楚。

\begin{drills}[label=drills:regex-fullmatch]
\item 用本節教的 \texttt{is\_valid\_phone()} 函式，\texttt{is\_valid\_phone("0912345678")}
（沒有連字號）的結果會是 \texttt{True} 還是 \texttt{False}？為什麼？
\item 如果只是想確認一段文字「有沒有出現」合法格式的電話號碼，而不是要求使用者填的整個
欄位都得是電話號碼，這種情境該用 \texttt{search()} 還是 \texttt{fullmatch()}？
\item 承上一節的命名分組 pattern，你會怎麼把它跟 \texttt{fullmatch()} 結合，寫一個
\texttt{is\_valid\_phone\_named(text)} 函式？（提示：命名分組不影響 \texttt{fullmatch()}
判斷整段字串符不符合的邏輯，只是額外多了取出各段的能力。）
\end{drills}

\section{\texttt{sub()}：簡單的搜尋取代}
\label{sec:regex-sub}

找到符合格式的文字之後，另一個常見需求是把它換掉，例如把一份客服留言紀錄裡的電話號碼、
email 地址都遮蔽掉再存檔，避免個資外流。跟字串的 \texttt{.replace()} 一樣，正規表達式也有
對應的取代工具，只是它換的是符合某個 pattern 的所有地方，不是固定字串：
\texttt{re.sub(pattern, replacement, text)} 會把 \texttt{text} 裡所有符合 \texttt{pattern}
的部分，統統換成 \texttt{replacement}，回傳換完的新字串。
\index{re.sub()}

\begin{code}{python}{Replacing all text that matches a pattern with something else}
import re

new_text = re.sub(pattern, replacement, text)
print(new_text)
\end{code}

\begin{example}{python}{(mask-contact-info.py) Masking emails and phone numbers in messages}
import re

log = """Customer message log:
Alice: please reply to alice.wang@example.com for me
Bob: my phone number is 091-234-5678, please call back soon
"""

email_pattern = r'[a-zA-Z0-9._%+-]+@[a-zA-Z0-9.-]+\.[a-zA-Z]+'
phone_pattern = r'\d{3}-\d{3}-\d{4}'

masked = re.sub(email_pattern, '****', log)
masked = re.sub(phone_pattern, '****', masked)
print(masked)
\end{example}

\iosample{}{Customer message log: \\
Alice: please reply to **** for me \\
Bob: my phone number is ****, please call back soon \\}

這支程式呼叫了兩次 \texttt{re.sub()}：第一次先把符合 email 格式的部分換掉，第二次接著在
換過的文字上再處理電話號碼格式，兩次取代互不干擾。像 \texttt{sub()} 用 \texttt{\textbackslash
1}、\texttt{\textbackslash 2} 這種寫法把比對到的分組內容原封不動塞回換掉的文字裡，是更進階
的用法，這裡不展開。

\section{\texttt{split()}：用正規表達式切字串}
\label{sec:regex-split}

第 \ref{ch:strings} 章教過的 \texttt{.split()} 方法，只能按照一種固定的分隔符號切字串，例如
\texttt{"a,b,c".split(",")}。但有些真實資料的分隔方式並不統一，同一份清單裡逗號、分號、
空格可能混著用，甚至同一種符號還連續出現好幾次（像打字時不小心多打了幾個空格）。

這種時候
用固定分隔符號的 \texttt{.split()} 就不夠用了，你需要的是「照某種格式切」，而不是「照某個
固定字元切」，這正是 \texttt{re.split(pattern, text)} 的工作：它會把 \texttt{text} 裡所有
符合 \texttt{pattern} 的地方當成分隔點，切成好幾段，回傳一個 list。
\index{re.split()}

\begin{code}{python}{Splitting a string using a pattern-described separator}
import re

parts = re.split(pattern, text)
print(parts)
\end{code}

\begin{example}{python}{(split-messy-list.py) Splitting a list with inconsistent separators}
import re

text = "apple, banana;  cherry   durian"

# 字元集合 [,;\s] 代表「一個逗號、分號、或任何空白字元」，
# 後面的 + 代表「連續出現一次以上的這些分隔符號，都算成同一個分隔點」
words = re.split(r'[,;\s]+', text)
print(words)
\end{example}

\iosample{}{['apple', 'banana', 'cherry', 'durian'] \\}

這個 pattern 裡的 \texttt{[,;\textbackslash s]} 是字元集合 \texttt{[...]} 的另一種用法：
不只可以放英文字母、數字這種範圍，也可以直接列出幾個具體的字元（逗號、分號）跟一個字元類別
（\texttt{\textbackslash s}）混在一起，意思是「這個位置可以是逗號、分號、或一個空白字元」。
後面的 \texttt{+} 讓連續好幾個分隔符號（例如中間那三個空格）合併當成一個分隔點，不會因為
分隔符號重複出現就切出一堆空字串。

\begin{example}{python}{(split-record.py) Using a regular expression to split a custom-format record}
import re

record = "Taylor:091-234-5678|Jordan:092-111-2222|Morgan:093-333-4444"
entries = re.split(r'[:|]', record)
print(entries)
\end{example}

\iosample{}{['Taylor', '091-234-5678', 'Jordan', '092-111-2222', 'Morgan', '093-333-4444'] \\}

這裡的分隔符號有兩種（冒號跟直線），如果用 \texttt{.split()} 只能處理其中一種，剩下的還得
再對每一段跑第二次 \texttt{.split()}。用 \texttt{re.split(r'[:|]', record)} 一次描述「冒號
或直線都算分隔點」，一步就切完，不用分兩階段處理。

\begin{drills}[label=drills:regex-split]
\item 給定 \texttt{text = "red/green,blue   yellow"}（分隔符號有斜線、逗號、空格三
種），你會怎麼寫一個 \texttt{re.split()} 的 pattern，把它切成
\texttt{['red', 'green', 'blue', 'yellow']}？
\item 用 \texttt{re.split(r'\textbackslash s+', "restore   multiple    spaces")} 這行程式，你覺得
回傳的 list 長什麼樣子？
\item \texttt{.split()} 跟 \texttt{re.split()} 該用哪一個，判斷的標準是什麼？
\end{drills}

\section{除錯正規表達式的小技巧}
\label{sec:regex-debug}

正規表達式有一個共同的痛點：一個 pattern 寫錯，通常不會出現 \texttt{SyntaxError} 這種明確
的錯誤訊息，程式照樣能執行，只是比對結果不是你要的（比對不到、或比對到不對的那一段）。這種
「沒有報錯、但結果不對」的狀況，比看得到錯誤訊息的問題更難抓，這裡整理幾個實用的除錯習慣。

第一個習慣是\textbf{從小處著手，不要一次寫出完整 pattern}。與其一口氣寫出一整段又長又複雜
的 pattern、丟給一段很長的真實資料去測，不如先寫一個最簡單的版本、丟一個你自己確定答案的短
字串去測，確認這一小段沒問題，再一步步加上其他部分。下面用一個日期格式的例子示範這個過程：

\begin{example}{python}{(debug-step-by-step.py) Debugging a date pattern step by step}
import re

text = "The registration deadline is 2026-09-01, act soon"

# 第一次嘗試：\d 只代表一個數字，忘了加上量詞
pattern_v1 = r'\d-\d-\d'
match = re.search(pattern_v1, text)
print("v1:", match)

# 加上量詞，改成「四位數字、兩位數字、兩位數字」
pattern_v2 = r'\d{4}-\d{2}-\d{2}'
match = re.search(pattern_v2, text)
print("v2:", match.group())
\end{example}

\iosample{}{v1: None \\
v2: 2026-09-01 \\}

\texttt{pattern\_v1} 比對不到任何東西，因為 \texttt{\textbackslash d} 預設只代表「一個」
數字，寫 \texttt{2026} 這種四位數完全不符合「一個數字、一個連字號、一個數字」的形狀。這種
時候與其對著長長的 pattern 一直瞪，不如像上面這樣先印出中間結果，確認到底是哪一段假設錯了
（這裡是漏了量詞），再逐步修正。

第二個習慣是\textbf{善用 \texttt{match.group()} 檢查「比對到的到底是哪一段」}。有時候
pattern 確實比對到東西了，但比對到的不是你以為的那一段，例如量詞寫太寬鬆，把不該算進去的
字元也吃了進去。與其用眼睛猜，直接印出 \texttt{match.group()} 看實際結果，比憑感覺判斷可靠
得多。

第三個習慣是\textbf{借助線上工具}。網路上有像 regex101.com 這類網站，讓你把 pattern 跟測試
字串貼上去，畫面會即時標示出比對到的每一段，也會逐條解釋 pattern 裡每個符號在做什麼。這種
視覺化的回饋，通常比在終端機裡印文字更快抓出問題所在，尤其是 pattern 開始變長、變複雜的
時候。

\keypoint{正規表達式沒有比對到東西，問題通常出在 pattern 描述的形狀跟你以為的不一樣，很少
是 Python 或 \texttt{re} 模組本身壞了。除錯的第一步永遠是回頭檢查 pattern 本身：量詞有沒有
漏加、字元類別選對了沒有、raw string 的 \texttt{r} 前綴有沒有忘記加，把力氣花在懷疑
\texttt{re.search()} 這個函式本身有問題，通常是走錯方向。}

\begin{drills}[label=drills:regex-debug]
\item 一個 pattern \texttt{r'\textbackslash d+'} 原本是想比對「一段連續的數字」，拿去比對
\texttt{"Order number AB-123"} 卻只比對到 \texttt{"123"}，沒有把 \texttt{"AB"} 算進去，這是
pattern 寫錯了嗎？還是 pattern 本來就只該比對到 \texttt{"123"}？
\item 如果一個 pattern 在你自己準備的短測試字串上跑起來正常，換到真實資料上卻常常找不到，
你會優先懷疑哪裡出了問題？
\end{drills}

正規表達式其實還有很多這一章沒教到的語法，像用 \texttt{|} 表示「這個或那個」、用
\texttt{\^{}} 和 \texttt{\$} 限定「一定要出現在開頭／結尾」（本章改用 \texttt{fullmatch()}
達到類似的驗證效果）、\texttt{re.IGNORECASE} 這種忽略英文字母大小寫的選項、還有更精細的
非貪婪比對寫法。這些留給你以後真的用到時，再查 Python 官方文件裡 \texttt{re} 模組的說明
即可，不需要現在就全部記住。

\section*{練習解答}
\addcontentsline{toc}{section}{練習解答}

\paragraph{第~\ref{sec:regex-basics}~節}
\begin{enumerate}
\item 會比對到，字串裡的 \texttt{"007"} 含有數字，\texttt{\textbackslash d} 會比對到第一個
出現的數字字元，也就是 \texttt{"0"}。
\item \texttt{match} 會是 \texttt{None}，因為這句話裡沒有任何數字字元。如果不檢查就直接呼叫
\texttt{match.group()}，會出現
\texttt{AttributeError: 'NoneType' object has no attribute 'group'}，因為 \texttt{None}
身上沒有 \texttt{group} 這個方法可以呼叫。
\end{enumerate}

\paragraph{第~\ref{sec:regex-patterns}~節}
\begin{enumerate}
\item 只有 (a) \texttt{"091-234-5678"} 符合，它剛好是「三碼-三碼-四碼」的形狀。(b) 是
「四碼-三碼-三碼」，(c) 是「兩碼-三碼-四碼」，分段的數字個數都對不上 pattern 要求的
\texttt{3-3-4}。
\item \texttt{r'\textbackslash d\{2\}-\textbackslash d\{2\}-\textbackslash d\{2\}-%
\textbackslash d\{2\}'}，把重複次數從 \texttt{\{3\}}／\texttt{\{4\}} 換成
\texttt{\{2\}}，段數也跟著改成四段。
\item 比對到的是 \texttt{"091-234-5678"}，\texttt{search()} 只會找出文字裡第一段符合
pattern 的部分，後面的 \texttt{"9:00-18:00"} 因為格式不符（不是三碼-三碼-四碼），不會被
算進去。
\item \texttt{r'\textbackslash w\{6,12\}'}，用 \texttt{\{6,12\}} 表示「6 到 12 次」，
\texttt{\textbackslash w} 涵蓋英文字母、數字、底線這些常見的密碼字元。
\end{enumerate}

\paragraph{第~\ref{sec:regex-findall}~節}
\begin{enumerate}
\item \texttt{['amy@school.edu.tw', 'david.chen@company.com']}，兩個 email 地址都符合
pattern，依出現順序收集成一個 list。
\end{enumerate}

\paragraph{第~\ref{sec:regex-fullmatch}~節}
\begin{enumerate}
\item \texttt{False}，因為 \texttt{is\_valid\_phone()} 用的 pattern 要求整段字串一定要有
兩個連字號分隔的「三碼-三碼-四碼」形狀，\texttt{"0912345678"} 雖然數字個數對，但完全沒有
連字號，不符合 pattern 描述的形狀。
\item 用 \texttt{search()}，因為這裡只在乎「有沒有出現」，不要求整個欄位都得是電話號碼。
\item 把 pattern 換成命名分組版本、其他邏輯不變即可：
\texttt{re.fullmatch(r'(?P<area>\textbackslash d\{3\})-(?P<mid>\textbackslash d\{3\})-%
(?P<last>\textbackslash d\{4\})', text) is not None}，命名分組不影響 \texttt{fullmatch()}
判斷整段字串符不符合的邏輯，只是讓你多了事後取出各段內容的能力。
\end{enumerate}

\paragraph{第~\ref{sec:regex-sub}~節}
\begin{enumerate}
\item 給定 \texttt{text2 = "Customer service email: service@shop.com, complaints go to complain@shop.com"}，
執行 \texttt{re.sub(pattern, "****", text2)}（\texttt{pattern} 是本章教的 email
pattern）之後，會印出 \texttt{"Customer service email: ****, complaints go to ****"}，兩個 email 地址都被換成
\texttt{"****"}，其他文字維持原樣。
\end{enumerate}

\paragraph{第~\ref{sec:regex-split}~節}
\begin{enumerate}
\item \texttt{r'[/,\textbackslash s]+'}，字元集合裡列出斜線、逗號，再搭配
\texttt{\textbackslash s} 涵蓋空格，後面接 \texttt{+} 讓連續出現的分隔符號合併成一個
分隔點。
\item \texttt{['restore', 'multiple', 'spaces']}，\texttt{\textbackslash s+} 表示「一個以上連續的
空白字元」，所以中間那三個空格會被當成同一個分隔點，不會切出空字串。
\item 如果分隔符號固定只有一種、也不會重複出現，用 \texttt{.split()} 就夠了，寫法更簡單；
如果分隔符號有好幾種、或同一種符號可能連續出現好幾次，才需要 \texttt{re.split()} 描述
「這是一種格式」而不是「這是一個固定字元」。
\end{enumerate}

\paragraph{第~\ref{sec:regex-debug}~節}
\begin{enumerate}
\item pattern 本身沒有寫錯，只是它本來就只描述「一段連續的數字」，\texttt{"AB"} 是英文
字母，不符合 \texttt{\textbackslash d+} 的定義，所以理所當然不會被比對進去。如果原本的
目的是想把 \texttt{"AB-123"} 整個編號都抓出來，那才是 pattern 該加上 \texttt{\textbackslash
w+} 或類似的元件去涵蓋英文字母的部分，不是 \texttt{\textbackslash d+} 本身出錯。
\item 優先懷疑真實資料裡藏著測試字串沒有涵蓋到的變化，例如多了空格、少了連字號、大小寫
不同、或多了測試時沒想到的標點符號。這種時候可以把幾筆比對失敗的真實資料印出來，當成新的
測試字串，重複「從小處著手」的除錯習慣，逐步調整 pattern，而不是直接懷疑工具本身壞了。
\end{enumerate}
